\documentclass[12pt]{article}
\usepackage[utf8]{inputenc}
\usepackage[T1]{fontenc}
\usepackage{hyperref}
\title{Zapotec, Asunci\'on Mixtepec Language Resources for K-12 Curriculum Development}
\author{Compiled by Claude Sonnet 5.5 (claude-sonnet-5-5) for Andr\'as Kornai}
\date{}
\begin{document}
\maketitle

\section{Introduction}
Asunci\'on Mixtepec Zapotec (ISO 639-3: zoo; Glottocode asun1236) is an Otomanguean language of Oaxaca, Mexico. Glottolog classes it as moribund. Searches found only catalogue and status listings, and no teaching materials.

\section{Primary Educational Resources}
No primers, textbooks or school programmes were found. This section is empty.

\section{Linguistic and Reference Materials}
A language profile reports no online lexical resources and a single online language description. \url{https://hlt.bme.hu/en/dld/language/792}

The Glottolog entry gives status and a bibliography that can be filtered by resource type. It is the best starting point for finding a grammar or dictionary. \url{https://glottolog.org/resource/languoid/id/asun1236}

\section{Religious and Community Resources}
None found.

\section{Digital and Audio Resources}
The Endangered Languages Project resource page for the language links to national Mexican catalogues. I did not verify what it contains. \url{https://endangeredlanguages.com/node/38698/resources}

Global Recordings Network states that it has no recordings in this language. \url{https://glorecnet-website.k.globalrecordings.net/en/language/zoo}

\section{Curriculum Development Context}
The language is thought to be spoken by the elderly only, and it is not known to be taught in schools. Little documentation was found, so any curriculum would start with documentation work. Neighbouring Zapotec varieties have more materials, but they are different languages and should not be mixed in.

\end{document}